% Options for packages loaded elsewhere
\PassOptionsToPackage{unicode}{hyperref}
\PassOptionsToPackage{hyphens}{url}
\documentclass[
  ignorenonframetext,
]{beamer}
\newif\ifbibliography
\usepackage{pgfpages}
\setbeamertemplate{caption}[numbered]
\setbeamertemplate{caption label separator}{: }
\setbeamercolor{caption name}{fg=normal text.fg}
\beamertemplatenavigationsymbolsempty
% remove section numbering
\setbeamertemplate{part page}{
  \centering
  \begin{beamercolorbox}[sep=16pt,center]{part title}
    \usebeamerfont{part title}\insertpart\par
  \end{beamercolorbox}
}
\setbeamertemplate{section page}{
  \centering
  \begin{beamercolorbox}[sep=12pt,center]{section title}
    \usebeamerfont{section title}\insertsection\par
  \end{beamercolorbox}
}
\setbeamertemplate{subsection page}{
  \centering
  \begin{beamercolorbox}[sep=8pt,center]{subsection title}
    \usebeamerfont{subsection title}\insertsubsection\par
  \end{beamercolorbox}
}
% Prevent slide breaks in the middle of a paragraph
\widowpenalties 1 10000
\raggedbottom
\AtBeginPart{
  \frame{\partpage}
}
\AtBeginSection{
  \ifbibliography
  \else
    \frame{\sectionpage}
  \fi
}
\AtBeginSubsection{
  \frame{\subsectionpage}
}
\usepackage{iftex}
\ifPDFTeX
  \usepackage[T1]{fontenc}
  \usepackage[utf8]{inputenc}
  \usepackage{textcomp} % provide euro and other symbols
\else % if luatex or xetex
  \usepackage{unicode-math} % this also loads fontspec
  \defaultfontfeatures{Scale=MatchLowercase}
  \defaultfontfeatures[\rmfamily]{Ligatures=TeX,Scale=1}
\fi
\usepackage{lmodern}
\ifPDFTeX\else
  % xetex/luatex font selection
\fi
% Use upquote if available, for straight quotes in verbatim environments
\IfFileExists{upquote.sty}{\usepackage{upquote}}{}
\IfFileExists{microtype.sty}{% use microtype if available
  \usepackage[]{microtype}
  \UseMicrotypeSet[protrusion]{basicmath} % disable protrusion for tt fonts
}{}
\makeatletter
\@ifundefined{KOMAClassName}{% if non-KOMA class
  \IfFileExists{parskip.sty}{%
    \usepackage{parskip}
  }{% else
    \setlength{\parindent}{0pt}
    \setlength{\parskip}{6pt plus 2pt minus 1pt}}
}{% if KOMA class
  \KOMAoptions{parskip=half}}
\makeatother
\setlength{\emergencystretch}{3em} % prevent overfull lines
\providecommand{\tightlist}{%
  \setlength{\itemsep}{0pt}\setlength{\parskip}{0pt}}
\usepackage{bookmark}
\IfFileExists{xurl.sty}{\usepackage{xurl}}{} % add URL line breaks if available
\urlstyle{same}
\hypersetup{
  hidelinks,
  pdfcreator={LaTeX via pandoc}}

\author{\texorpdfstring{}{}}
\date{}

\begin{document}

\section{IR progression, translation engine, and
algorithms}\label{ir-progression-translation-engine-and-algorithms}

\begin{frame}[fragile]{Progressive IRs}
\protect\phantomsection\label{progressive-irs}
Processing proceeds through a sequence of representations. Which stages
are complete for a given repository is summarized in
\href{../cogant/docs/reference/implementation_status.md}{\texttt{../cogant/docs/reference/implementation\_status.md}}
(partial areas include translation rules, state space, and Rust
acceleration):

\begin{enumerate}
\tightlist
\item
  \textbf{Repo IR} --- entities and relationships extracted from
  parsers.
\item
  \textbf{Program graph IR} --- consolidated graph with deduplication
  and metadata.
\item
  \textbf{Semantic mapping IR} --- output of the translation rule
  engine.
\item
  \textbf{State space IR} --- variables, actions, transitions,
  observations.
\item
  \textbf{Process model IR} --- higher-level control patterns where
  implemented.
\item
  \textbf{Validation IR} --- coverage, confidence analysis, schema
  checks.
\end{enumerate}

COGANT's program graph sits in the same conceptual space as compiler
intermediate representations such as LLVM {[}@lattner2004llvm{]} and
MLIR {[}@lattner2021mlir{]}, but targets behavioral extraction and
export for downstream learning rather than code generation or
optimization.
\end{frame}

\begin{frame}[fragile]{Translation rules}
\protect\phantomsection\label{translation-rules}
The \textbf{translation} stage applies declarative rules that refine
roles, attach labels, and adjust confidence. Concurrency targets and
layering are described in
\texttt{../cogant/docs/architecture/README.md}. The rule engine composes
passes over the graph in a \textbf{fixpoint loop}: rules are re-applied
until no new semantic mappings emerge or a configurable iteration cap is
reached. The fixpoint loop follows the classical formulation of program
analysis as fixpoint computation over a lattice of abstract states
{[}@cousot1977abstract{]}; in our setting the lattice is the set of
semantic mappings partially ordered by inclusion, and the monotone
operator is the composition of all registered rule applications. In the
shipped implementation, each pass applies rules in descending
\texttt{rule.priority}, and conflict resolution later compares
\texttt{(rule\_priority,\ confidence\_score)} tuples when two mappings
overlap, following the principle that edit representations should be
composable {[}@yin2019learning{]}.

\begin{block}{Fixpoint iteration, conflict resolution, and coverage}
\protect\phantomsection\label{fixpoint-iteration-conflict-resolution-and-coverage}
The shipped \texttt{TranslationEngine} in
\texttt{../cogant/py/cogant/translate/engine.py} realizes the fixpoint
loop concretely. Each iteration walks every registered rule once: the
engine calls \texttt{rule.matches(graph,\ query)} to collect candidate
fragments, then calls \texttt{rule.apply(graph,\ match)} on each,
accumulating any resulting \texttt{SemanticMapping} objects keyed by
their stable IDs. A per-pass counter tracks mappings that were genuinely
new (not already present in the running set), and the loop terminates as
soon as a pass completes with zero additions. The engine logs each
iteration boundary through an internal match log, so the number of
passes required to reach a fixed point is directly observable in the
post-run diagnostics. The default iteration cap is
\texttt{max\_iterations\ =\ 10}; in testing, most repositories converge
well before that bound, and the cap exists primarily as a safety valve
against pathological rule sets that could otherwise oscillate
indefinitely.

After fixpoint termination, the engine invokes
\texttt{\_resolve\_conflicts()} to reconcile mappings whose
\texttt{graph\_fragment\_node\_ids} sets overlap. For each overlapping
pair the engine retains the mapping with the larger
\texttt{(rule\_priority,\ confidence\_score)} key and discards the
other, logging a \texttt{conflict\_resolved} event that records the
losing ID, the winning ID, and the specific overlap set. Most shipped
rules use the default \texttt{rule.priority} of 0; the
mutating-subsystem / hidden-state rule uses priority 1 so it survives
overlaps with same-confidence class-level aggregates (for example
containment summaries) on the same \texttt{CLASS} node. A companion
entry point, \texttt{translate\_with\_confidence()}, runs the standard
fixpoint loop, rescores every surviving mapping through the
\texttt{ConfidenceModel}, and then re-resolves conflicts so that any
ordering shifts induced by rescoring are honored.

Translation coverage -- the fraction of graph nodes that received at
least one semantic mapping -- is reported by
\texttt{get\_coverage\_report(graph)}. It returns the total node count,
the number of covered nodes, the number of uncovered nodes, a
\texttt{coverage\_percent} value rounded to two decimal places, and the
sorted list of uncovered node IDs. The uncovered list is intentionally
emitted verbatim so that downstream tooling can target unmapped regions
for manual review or rule extension, rather than burying the gap behind
a single aggregate number {[}@allamanis2018survey{]}.

\begin{algorithm}
\caption{Fixpoint translation engine}
\label{alg:fixpoint}
\KwIn{Program graph $G = (V, E)$; rule set $R$; maximum iterations $K$ (default 10)}
\KwOut{Set of semantic mappings $\mathcal{M}$ keyed by stable ID}
$\mathcal{M} \leftarrow \emptyset$\;
\For{$k \leftarrow 1$ \KwTo $K$}{
    $n_{\text{new}} \leftarrow 0$\;
    \ForEach{rule $r \in R$ sorted by $\text{priority}(r)$ descending}{
        \ForEach{match $m \in r.\text{matches}(G)$}{
            $\mu \leftarrow r.\text{apply}(G, m)$\;
            \If{$\mu \neq \bot$ \textbf{and} $\mu.\text{id} \notin \mathcal{M}$}{
                $\mathcal{M} \leftarrow \mathcal{M} \cup \{\mu\}$\;
                $n_{\text{new}} \leftarrow n_{\text{new}} + 1$\;
            }
        }
    }
    \If{$n_{\text{new}} = 0$}{
        \textbf{break} \tcp*{fixed point reached}
    }
}
$\mathcal{M} \leftarrow \textsc{ResolveConflicts}(\mathcal{M})$\;
\Return $\mathcal{M}$\;
\end{algorithm}

Algorithm \ref{alg:fixpoint} summarizes the engine. Termination is
guaranteed because each iteration either produces at least one new
mapping (whose stable ID is then fixed in \(\mathcal{M}\)) or terminates
by the break condition; the outer \(K\) bound serves as a safety valve
for pathological rule sets.

\begin{algorithm}
\caption{Priority-ordered conflict resolution}
\label{alg:conflict}
\KwIn{Mapping set $\mathcal{M}$; priority function $p(\cdot)$}
\KwOut{Reduced mapping set with no overlapping fragments}
Build inverted index $\mathcal{I}: V \to 2^{\mathcal{M}}$ from node IDs to mappings touching them\;
$\mathcal{C} \leftarrow \emptyset$ \tcp*{conflict pairs}
\ForEach{node $v$ with $|\mathcal{I}(v)| \geq 2$}{
    \ForEach{$(\mu_a, \mu_b) \in \binom{\mathcal{I}(v)}{2}$}{
        $\mathcal{C} \leftarrow \mathcal{C} \cup \{(\mu_a, \mu_b)\}$\;
    }
}
$\mathcal{R} \leftarrow \emptyset$ \tcp*{removal set}
\ForEach{$(\mu_a, \mu_b) \in \mathcal{C}$}{
    \If{$\mu_a \in \mathcal{R}$ \textbf{or} $\mu_b \in \mathcal{R}$}{\textbf{continue}}
    $k_a \leftarrow (p(\mu_a), c(\mu_a))$; $k_b \leftarrow (p(\mu_b), c(\mu_b))$\;
    \eIf{$k_a \geq k_b$}{$\mathcal{R} \leftarrow \mathcal{R} \cup \{\mu_b\}$}{$\mathcal{R} \leftarrow \mathcal{R} \cup \{\mu_a\}$}
}
\Return $\mathcal{M} \setminus \mathcal{R}$\;
\end{algorithm}

Algorithm \ref{alg:conflict} detects conflicts via an inverted index in
\(O(\sum_v |\mathcal{I}(v)|^2)\) worst-case time, which is substantially
faster than the naive all-pairs scan for graphs where most nodes carry
at most one mapping.
\end{block}
\end{frame}

\end{document}
